\section{Related work}\label{related-work}

\textbf{World models and world generators.} Genie learns an interactive environment from video \citep{bruce2024genie}. WorldGen generates traversable, editable worlds from text \citep{wang2025worldgen}; WorldAct decomposes a monolithic generated world into objects after the fact \citep{hu2026worldact}; WorldSculpt composes per-object meshes from grounded video \citep{niu2026worldsculpt}. SCORE ships no world-model output; a world generator may appear only as an optional reference step (Section 3.6).

\textbf{Agents that build scenes.} Holodeck has a language model write spatial constraints for a solver \citep{yang2024holodeck}; SceneCraft and recursive code world models write scenes as programs and compare renders with a reference \citep{hu2024scenecraft, li2026rcwm}; WorldClaw and AutoUE use agents to assemble open worlds and game code \citep{guo2026worldclaw, yin2026autoue}; 3D-RE-GEN and SceneConductor decompose one picture into placed objects \citep{sautter2025regen, kim2026sceneconductor}. LEGO-Anything is closest in spirit: its coding agents show ``weak scene initialization, regressive edits during iteration, and unreliable self-evaluation'', addressed with tools rather than training \citep{li2026lego}. SCORE shares that position but builds worlds that do not exist yet, to a person's picks.

\textbf{Procedural generation, parts and style.} Infinigen and Infinigen Indoors generate worlds and rooms from rules, with materials as their own generators \citep{raistrick2023infinigen, raistrick2024indoors}; ProcFunc gives these generators an interface on which language models make far fewer coding errors \citep{raistrick2026procfunc}. PartCrafter and Point2Part split shapes into parts \citep{lin2025partcrafter, tsui2026point2part}. Keeping one style across generated assets is named as open \citep{wu2026production, yang2026flowscene}; Hunyuan3D Studio fixes it at the picture \citep{lei2025hunyuanstudio}. SCORE instead takes surface out of the generated models altogether.
